\scanpage{83}
\mspage{9}
by a torus --- in practice, by a suitable ``generalized Jacobian'' of $X$.

But one must remember that for such a variety (for example an elliptic curve, or the variety $G_m(\mathbb C)=\mathbb C^*$) the action of $\Gamma$ on the groupoid $\Pi_1(X_K)$ is no longer at all faithful, so to ``identify'' a $u\in\Gamma$, it is not enough to know invariants of the type $g_l(u)$ \uline{relative to} \uline{this same} $X$. It will therefore be essential to introduce ``identification parameters'' for $u$, using other algebraic varieties defined over $K$ --- the most natural being
\[
\tag{35}\unsure{\mathbb U|U}_{0,3}=\mathbb P^1\setminus\{0,1,\infty\}
\]
\figure{\begin{tikzpicture}\draw(0,0)--(2.8,.12);\end{tikzpicture}}
Now letting $X$ vary, by \unsure{considering} a morphism of $K$-schemes
\[
\tag{36}f:X\unsure{\mathbin{-}|\to}X
\]
one deduces morphisms of groupoids
\margin{82}
